% ------------------------------------------------------------------------
% -*-TeX-*- -*-Hard-*- Smart Wrapping
% ------------------------------------------------------------------------


\chapter{考察}
\label{cha:5}

\section{実験結果の解釈: 議論の効果は2軸のトレードオフで決まる}
\label{sec:5.1}

実験1の最も重要な発見は，マルチエージェント議論の効果が単一の符号を持たず，ベンチマークによって系統的に反転することである．素の議論はMATH-500で$+4.0$pt・SuperGPQAで$+3.1$ptと有意に有効だが，MMLU-Proでは$-4.1$ptと有意に有害だった．逆にLoRAペルソナチームはMMLU-Proでのみ有効だった．この対照は，議論チームの性能を（i）\textbf{多様性の利得}（異なる推論経路が正解を発見する確率の向上）と（ii）\textbf{個体能力の毀損}（ペルソナ注入が支払うコスト）という2軸のトレードオフとして整理すると一貫して説明できる．

MMLU-Proのような知識推論では，温度サンプリングだけでは推論経路の多様性が乏しく（実験1で素の議論はベース単体より悪化），重みに焼き込まれたペルソナが供給する視点の多様性が利得をもたらす．一方でペルソナSFTの能力毀損は比較的軽微であり，差し引きで正となる．MATH-500やSuperGPQAのような長い導出を要する領域では，多様性は温度サンプリングでも十分に確保できる（素の議論やSC@9が強い）一方，SFTによる能力毀損が致命的となり，差し引きが負に転じる．

この整理は先行研究間の見かけ上の矛盾も解消する．「議論は多数決に劣る」とする理論的・実証的報告と「議論が小型モデルで有効」とする報告は，領域の検証可能性と能力律速性の違いとして両立する．特にMATH-500における素の議論の有意な効果は，「議論単体は信念のマルチンゲールを誘導し期待正解率を改善しない」という近年の定理的整理\cite{choi2025debate}に対する\textbf{領域依存の反例}を与える．検証可能な領域では，他者の解の正誤を自己検算によって判定できるため信念更新が訂正方向へ偏り，マルチンゲール性（更新の無情報性）の前提が破れる，という解釈が可能である．実測でもMATH-500の集約損失（oracle上限と実チーム性能の差）はほぼゼロであり（表\ref{tab:4.oracle}），検証可能性が集約の質を規定することを支持する．

\section{損失分解が導いた処方とその効果}
\label{sec:5.2}

実験1の診断（\S\ref{sec:4.6.2}）はチームの誤差を個体能力損失と集約損失に分解し，それぞれに対応する処方を実験2で検証した．

\textbf{能力毀損の機構と修復}．ペルソナSFTがベースモデルの思考連鎖を半分以下（中央値2{,}423字$\rightarrow$1{,}188字）に圧縮していたという生成ログの観察は，毀損の主因が「短い応答スタイルの刷り込みによる長い導出能力の選択的破壊」であることを示す．これはSFTデータとモデル分布のギャップが忘却を生むという近年の知見と整合する．処方としてベースモデル自身が生成した検証済みの長い思考連鎖（平均6{,}503字）を約38\%混合して再学習したところ，solo精度は3ベンチマークすべてで基準を満たし，一部はベース単体を上回った（G1）．重要なのは，この修復が\textbf{ペルソナの多様性を保ったまま}達成された点である（同一測定環境で再学習チームは旧チームをMMLU-Pro $+3.7$pt・MATH-500 $+6.3$pt上回っており，ペルソナ間の視点差は失われていない）．ペルソナSFTは「能力を代償に多様性を買う取引」だが，リプレイはこの取引のパレート境界を押し広げる．

\textbf{集約損失と検証型集約}．SuperGPQAでは3体のうち少なくとも1体が正解している問題が51.7\%あるにもかかわらず，多数決＋議論は38.7\%しか回収していなかった．logprobに基づく重み付き投票は$+0.7$ptの改善に留まり，比較選択型の裁定（GenSelect型）は上積みを示さなかった（G2）．集約損失の大部分は依然として未回収であり，「正解が少数派である問題」の回収は4B級では未解決の課題として残る．なお実験2でsolo能力が修復された結果，議論後に票が割れる問題自体が減少したため，集約改善の寄与余地が実験1の見積もりより縮小した点には注意が必要である．

\textbf{条件付き更新と匿名化}．「自己の推論に具体的な誤りを特定できた場合のみ回答変更を許す」条件付き更新はパイロットで$+1.0$pt（非有意）と小さいが悪化のない効果を示し，発言の匿名化と併せて採用した．これらは追従（sycophancy）による正解の破壊——実験1では500問中26問で正しい多数派が議論後に崩壊していた——への対処であり，議論の信念更新を訂正方向へ偏らせる介入として位置づけられる．

\section{進化的最適化の再解釈: 選抜ではなく修復}
\label{sec:5.3}

本研究の当初の中心仮説は「チームレベル適応度（厳密Shapley値）による選択圧が協調的な個体を選抜する」であった．実験1の結果はこの仮説を部分的にしか支持しない．進化の寄与はMATH-500の$+4.3$pt（有意）のみで，選抜対象としたMMLU-Proではゼロだった．機構分析はより根本的な再解釈を与える．進化前後でLoRAの$\Delta W$のFrobeniusノルムは全役割で完全に一致（比1.00）しており，進化は重みの大きさではなく\textbf{方向のみを回転}させていた．交叉（重みブレンド）の平均化作用がSFTで注入された損傷成分を打ち消した，すなわち\textbf{進化は選抜機構としてではなく修復機構として機能した}と解釈できる．

この解釈は2つの含意を持つ．第一に，適応度100問の標準誤差$\pm 4.9$ptの下で毎世代最大値選抜を行う設計は，1〜2pt差の真の優劣を判別できず（必要問題数の試算は1判断あたり約2{,}600問），開発セット上の見かけの改善（$+12$pt）は勝者の呪いとして説明される．チームレベル適応度の理論的魅力を実証するには，識別力の高い問題の厳選や逐次半減法などによる選抜ノイズの抑制が前提条件となる．第二に，修復が目的ならば進化的探索は迂遠であり，実験2のリプレイ再学習が同じ役割をより安価かつ確実に果たした（再学習チームは進化後チームをプールで$+3.2$pt上回る）．本研究では実験2でこの上位互換性が確認されたため，改良版進化ループの実行よりも処方の検証を優先した．

\section{Self-Consistencyとの比較の最終形}
\label{sec:5.4}

計算量を揃えたSelf-Consistency（SC@9）は本研究を通じて最強のベースラインであり続け，測定環境を統一した最終評価では，診断に基づく処方を尽くした再学習チームですらプールで$-3.0$pt（$p<10^{-6}$）と有意に及ばなかった．「評価設計を厳密にするとMADの優位は消える」という批判\cite{smit2024should}を，（i）計算量マッチ，（ii）対応あり検定と多重比較補正，（iii）能力毀損の修復と集約改善という対策の実施，（iv）測定環境の統制，のすべてを満たした上で確認した結果であり，本研究で最も統制の効いた知見である．

敗因の構造は明快である．SC@9の強さは「同一分布からの独立な9サンプルの多数決」がもたらす分散削減にあり，ベース単体（1サンプル）に対して$+3.1$pt（$p<10^{-6}$）を安定して稼ぐ．一方チーム議論は6生成をペルソナ3体$\times$2ラウンドに配分するため，独立サンプル数では3にすぎず，議論による情報交換がその差を埋める必要がある．実験2の集約改善後もこのギャップは埋まらなかった．例外はMATH-500であり，チームはSC@9と統計的同等（$-0.7$pt，$p=0.27$），ベース単体には$+3.8$pt（$p=0.0003$）と有意に優る．検証可能領域では議論の1ラウンドが誤答の訂正として機能し，少ないサンプル数を補償することを示す．

なお実験2の途中経過（測定環境の統制前）では「チームがSC@9をMMLU-Proで$+5.5$pt上回る」という結果が一時得られたが，これは評価イメージ間の系統差（\S\ref{sec:4.6.4}）の産物であった．系統差はベースモデルの素の生成に選択的に現れLoRA適用時にはほぼ現れないため，\textbf{ベース系条件だけが過小評価され，チーム側が見かけ上優位になる}という質の悪いバイアスを生む．マルチエージェント議論の評価においてベースライン測定の環境同一性がいかに重要かを示す実例である．

\section{本研究の限界}
\label{sec:5.5}

\begin{itemize}
\item \textbf{規模の限界}: 3エージェント・単一モデル系列（Qwen3-4B）・3ベンチマークでの結果であり，より大きなモデルや異種モデル混成チームへの一般化は未検証である．特に「能力毀損」はベースモデルの事後学習の質に依存するため，モデル系列が変われば定量関係は変わりうる．
\item \textbf{実験1と実験2の数値の非互換}: 評価イメージ間の系統差（\S\ref{sec:4.6.4}）により，実験1（旧環境）と実験2（新環境）の絶対値は直接比較できない．実験1の内部比較（領域依存性・損失分解）と実験2の内部比較（処方の効果・対SC@9）はそれぞれ同一環境内の対応あり比較として有効だが，例えば「実験1の素の議論と実験2の再学習チーム」のような環境跨ぎの比較は本論文では行っていない．
\item \textbf{シード数の非対称}: 最終評価はSC@9と再学習チームが6シード，ベース単体と旧チームが3シードであり，後者を含む検定は検出力がやや低い．
\item \textbf{進化ループの条件}: 実験1の進化は選抜ノイズが大きい設定（固定100問・最大値選抜）で実行されており，本研究の「進化=修復」という解釈はこの設定に条件づけられる．識別力厳選と逐次半減法を備えた進化（設計は研究設計書v3 \S 3.4として保存）が選抜として機能するかは未検証である．
\item \textbf{進化ループの生成条件}: 進化ループは最大512トークン・top-p 0.9の旧生成設定で実行されており，最終評価（8{,}192トークン・top-p 0.8）と生成条件が異なる．ループ内部の比較は全個体で条件が揃っているため公正だが，進化時と評価時の挙動差は残る．
\item \textbf{未実施のアブレーション}: solo適応度進化（A1）等は計算資源の制約から未実施であり，チームレベル適応度の寄与の直接分離は今後の課題である．
\item \textbf{評価領域}: 全ベンチマークが英語の選択式・短答式であり，生成型タスクや日本語タスクへの外的妥当性は限定的である．
\end{itemize}
